% SHARD-ID: AQM-33-FULL-DUAL-LOCAL-MODELS-FOR-KT
% SHARD-TITLE: Full-Dual and Local Models for \(k(t)\)
% SHARD-SUMMARY: Explains why the algebraic \(k(t)\) representation is not the full Stone-von Neumann theorem.
% SHARD-SUMMARY: Records the full discrete-dual, local-field, and adelic completions where analytic uniqueness belongs.
% SHARD-SUMMARY: Works explicit \(k=\mathbb F_3\) commutation examples.
% SHARD-KEYWORDS: k(t), Pontryagin dual, local field, adele, Stone-von Neumann, F3, Weyl-Heisenberg

\section{Full-Dual and Local Models for \texorpdfstring{\(k(t)\)}{k(t)}}
\label{sec:full-dual-local-models-for-kt}

\subsection{Status and Source Anchors}

This shard completes AQM-32.  The full-dual LCA statement uses Prasad,
\path{references/heisenberg_weil/prasad_0912_0574_source/SvN.tex}, lines
95--107 and 164--190, and Beny--Crann--Lee--Park--Youn,
\path{references/heisenberg_weil/beny_crann_lee_park_youn_2204_08162_source/main.tex},
lines 461--473.  The warning about countably infinite discrete rings uses
Bekka, \path{references/heisenberg_weil/bekka_2502_00387_source/CCR-GeneralRings-v5.tex},
line 1014.  Local and adelic background is registered in
\path{references/heisenberg_weil/SOURCES.md} from Poonen's Tate notes and
Kudla's local theta notes.

\subsection{Lamport Dependency Skeleton}

\[
\begin{array}{rcl}
D1&:&K_{\rm disc}\text{ is }K=k(t)\text{ with the discrete topology.}\\
L1&:&K_{\rm disc}\not\simeq\widehat{K_{\rm disc}}\text{ as LCA groups.}\\
D2&:&G_{\rm full}=K_{\rm disc}\times\widehat{K_{\rm disc}}\text{ is the full
      LCA Weyl phase space.}\\
L2&:&\text{The algebraic }K\text{-modulations form a dense subgroup of }
      \widehat{K_{\rm disc}}.\\
D3&:&K_\infty=k((t^{-1}))\text{ is the local completion at infinity.}\\
T1&:&\text{Stone-von Neumann uniqueness belongs to }G_{\rm full}\text{ or
      to local/adelic LCA completions.}\\
E1&:&k=\mathbb F_3\text{ gives explicit }\omega\text{-commutation examples.}
\end{array}
\]

\subsection{The Full Discrete-Dual Model D1--L2}

The algebraic representation of AQM-32 is an algebraic Schrödinger
representation theorem for the concrete group \(H_K\).  It is not the full LCA
Stone-von Neumann theorem.  The reason is structural: if \(K\) is given the
discrete topology, then its Pontryagin dual \(\widehat{K_{\rm disc}}\) is
compact.  Since \(K\) is countably infinite, \(K\) is not isomorphic to
\(\widehat{K_{\rm disc}}\) as an LCA group.  Bekka records the same obstruction
for countably infinite discrete rings.

The full LCA Weyl system with configuration group \(K_{\rm disc}\) is
\[
  G_{\rm full}=K_{\rm disc}\times\widehat{K_{\rm disc}}.
\]
It acts on \(\ell^2(K)\) by
\[
  W(a,\gamma)=T_aM_\gamma,\qquad
  M_\gamma|\xi\rangle=\gamma(\xi)|\xi\rangle.
\]
Prasad's theorem applies to this \(G_{\rm full}\).  The algebraic \(H_K\) is
the restriction to the subgroup of characters
\[
  \gamma_b(\xi)=\chi_K(b\xi),\qquad b\in K.
\]
By AQM-32 \(L1\), \(b\mapsto\gamma_b\) is injective.

\paragraph{Density derivation.}
Let \(\Gamma_K=\{\gamma_b:b\in K\}\subset\widehat{K_{\rm disc}}\).  Its
annihilator in the double dual \(K_{\rm disc}\) is
\[
  \Gamma_K^\perp
  =
  \{\xi\in K:\gamma_b(\xi)=1\text{ for all }b\in K\}.
\]
By nondegeneracy of \((b,\xi)\mapsto\chi_K(b\xi)\), this annihilator is
\(\{0\}\).  Pontryagin duality identifies the annihilator of the closure
\(\overline{\Gamma_K}\) with \(\Gamma_K^\perp\).  Hence
\(\overline{\Gamma_K}=\widehat{K_{\rm disc}}\).  The algebraic model therefore
uses a dense rational-function family of modulations, while the full LCA
theorem uses every continuous character of \(K_{\rm disc}\).

\subsection{Local and Adelic Completions D3--T1}

For analytic arithmetic quantum fields, the more natural locally compact
configuration spaces are local completions such as
\[
  K_\infty=k((t^{-1}))
\]
or, later, the adele ring \(\mathbb A_K\).  A nontrivial continuous additive
character identifies a local field with its Pontryagin dual, so the canonical
Weyl system has phase space \(K_\infty\times K_\infty\) after that
identification.  The Schrödinger representation then lives on
\[
  L^2(K_\infty).
\]
For the adelic theory, a residue character built from a differential gives the
standard self-dual adelic setting; \(K\) embeds diagonally as arithmetic input,
not as a replacement for the full dual phase group.

Thus there are three honest objects:
\[
\begin{array}{ll}
\ell^2(K): & \text{discrete algebraic rational-function representation;}\\
\ell^2(K_{\rm disc})\text{ with }\widehat K: &
  \text{full LCA discrete-dual Stone-von Neumann system;}\\
L^2(K_\infty)\text{ or }L^2(\mathbb A_K): &
  \text{local or global analytic arithmetic system.}
\end{array}
\]
They should not be collapsed into one notation.

\subsection{Concrete Examples for \texorpdfstring{\(k=\mathbb F_3\)}{k=F3} E1}

Let \(k=\mathbb F_3\), \(K=\mathbb F_3(t)\), \(\tau=\operatorname{id}\), and
\[
  \omega=\exp(2\pi i/3).
\]
Then
\[
  \chi_K(f)=\omega^{[t^{-1}]f}.
\]
For \(a=1\) and \(b=t^{-1}\),
\[
  M_{t^{-1}}T_1=\omega\,T_1M_{t^{-1}}.
\]
Indeed, \(ba=t^{-1}\), whose \(t^{-1}\)-coefficient is \(1\).

For \(a=t\) and \(b=t^{-1}\),
\[
  M_{t^{-1}}T_t=T_tM_{t^{-1}},
\]
because \(ba=1\) has \(t^{-1}\)-coefficient \(0\).  For \(a=t\) and
\(b=t^{-2}\),
\[
  M_{t^{-2}}T_t=\omega\,T_tM_{t^{-2}},
\]
because \(ba=t^{-1}\).

On basis vectors,
\[
  T_{t+1}|t^{-1}\rangle=|t+1+t^{-1}\rangle,
\]
\[
  M_{t^{-1}}|1\rangle=\omega|1\rangle,\qquad
  M_{t^{-1}}|t\rangle=|t\rangle,
\]
and
\[
  M_{t^{-2}}|t\rangle=\omega|t\rangle.
\]
Thus the rational-function Heisenberg group already contains infinitely many
explicit qutrit-like commutation relations, indexed by rational functions
rather than by closed points.  This is a different layer from the
closed-point quasi-local field of \(\Spec\mathbb F_3[t]\), where each monic
irreducible \(\pi\) contributes one finite qudit
\(\ell^2(\mathbb F_3[t]/(\pi))\).

\subsection{Conclusion}

The concrete rational-function construction gives an irreducible unitary
representation on \(\ell^2(k(t))\).  Stone-von Neumann uniqueness belongs to a
larger or completed LCA object: the full discrete-dual system, a local
completion, or an adelic completion.  The next arithmetic-field step should
choose which completion is physically intended before importing Fourier
transform, Gaussian, or Weil-representation statements.
